% Reader-facing teaching elements. Every chapter uses the same small set so a
% reader learns the visual grammar once:
%   \ChapterGoals{<objectives>}{<prerequisites>}  chapter opening
%   \begin{KeyIdeas} ... \end{KeyIdeas}            chapter synthesis
%   \begin{BioBridge}  for readers arriving from computing
%   \begin{CSBridge}   for readers arriving from biology
%   \begin{Pitfall}    a misconception, failure case or negative control
%   \begin{Frontier}   current research, with explicit scope of the claim
%   \begin{TryThis}    predict before reading on
%   \begin{Worked}{title} a fully worked example
\usepackage[skins,breakable]{tcolorbox}
\tcbset{bookbase/.style={enhanced,breakable,sharp corners,boxrule=0pt,
  left=9pt,right=8pt,top=6pt,bottom=6pt,before skip=10pt,after skip=10pt,
  parbox=false}}

\newcommand{\BoxMark}[1]{\textcolor{#1}{\rule[-1pt]{7pt}{7pt}}\hspace{5pt}}
% \BoxHead{colour}{title}: the heading line placed at the top of each box.
\newcommand{\BoxHead}[2]{{\sffamily\bfseries\color{Ink}\BoxMark{#1}#2\par}\smallskip}

\newtcolorbox{KeyIdeas}{bookbase,colback=DNAblue!6,
  borderline west={3pt}{0pt}{DNAblue},
  before upper={\setlength{\parskip}{4pt}\raggedright\BoxHead{DNAblue}{Key ideas}}}
\newtcolorbox{BioBridge}{bookbase,colback=ProteinGold!12,
  borderline west={3pt}{0pt}{ProteinEdge},
  before upper={\setlength{\parskip}{4pt}\raggedright\BoxHead{ProteinEdge}{Biology bridge}}}
\newtcolorbox{CSBridge}{bookbase,colback=CompTeal!8,
  borderline west={3pt}{0pt}{CompTeal},
  before upper={\setlength{\parskip}{4pt}\raggedright\BoxHead{CompTeal}{Computing bridge}}}
\newtcolorbox{Pitfall}{bookbase,colback=RejectRed!6,
  borderline west={3pt}{0pt}{RejectRed},
  before upper={\setlength{\parskip}{4pt}\raggedright\BoxHead{RejectRed}{Common pitfall}}}
\newtcolorbox{Frontier}{bookbase,colback=EvidenceViolet!7,
  borderline west={3pt}{0pt}{EvidenceViolet},
  before upper={\setlength{\parskip}{4pt}\raggedright\BoxHead{EvidenceViolet}{Research frontier}}}
\newtcolorbox{TryThis}{bookbase,colback=StrandGreen!7,
  borderline west={3pt}{0pt}{StrandGreen},
  before upper={\setlength{\parskip}{4pt}\raggedright\BoxHead{StrandGreen}{Predict first}}}
\newtcolorbox{Worked}[1]{bookbase,colback=PanelBg,
  borderline west={3pt}{0pt}{InkSoft},
  before upper={\setlength{\parskip}{4pt}\raggedright\BoxHead{InkSoft}{Worked example: #1}}}

% Chapter opening: learning goals beside the prerequisites they rely on.
\newcommand{\ChapterGoals}[2]{%
 \begin{tcolorbox}[enhanced,breakable,sharp corners,boxrule=0pt,colback=PanelBg,
   borderline north={2pt}{0pt}{DNAblue},left=9pt,right=9pt,top=7pt,bottom=7pt,
   before skip=4pt,after skip=12pt,parbox=false,
   before upper={\setlength{\parskip}{3pt}\raggedright}]
 {\sffamily\bfseries\color{DNAblue}After this chapter you should be able to}\par
 \begin{itemize}\setlength{\itemsep}{1pt}\setlength{\parskip}{0pt}#1\end{itemize}
 \par\smallskip
 {\sffamily\bfseries\color{InkSoft}You will use}\quad #2
 \end{tcolorbox}}

% The original single-purpose callout keeps its interface, in the new style.
\renewcommand{\Callout}[2]{%
 \begin{tcolorbox}[bookbase,colback=DNAblue!6,borderline west={3pt}{0pt}{DNAblue},
   before upper={\setlength{\parskip}{4pt}\raggedright\BoxHead{DNAblue}{#1}}]#2\end{tcolorbox}}

% Part openers: a short orientation page for each group of chapters.
% \BookPart{title}{key}{orientation text, usually ending with a chapter list}
\newcommand{\BookPart}[3]{%
 \clearpage\thispagestyle{empty}%
 \refstepcounter{part}%
 \addcontentsline{toc}{part}{\protect\makebox[3.2em][l]{\thepart}#1}%
 \label{part:#2}%
 \noindent\begin{tikzpicture}[remember picture,overlay]
  \fill[DNAblue] (current page.north west) rectangle ([yshift=-3.6cm]current page.north east);
  \begin{scope}[shift={([xshift=1.2cm,yshift=-1.8cm]current page.north west)}]
   \foreach \hx in {0,0.22,...,18.6}{
    \draw[white,opacity=.35,line width=.9pt] (\hx,{.55*sin(360*\hx/3.4)}) -- (\hx,{.55*sin(360*\hx/3.4+150)});}
   \draw[white,line width=2.2pt,domain=0:18.6,samples=200,smooth] plot (\x,{.55*sin(360*\x/3.4)});
   \draw[DNAsky!60,line width=2.2pt,domain=0:18.6,samples=200,smooth] plot (\x,{.55*sin(360*\x/3.4+150)});
  \end{scope}
 \end{tikzpicture}\par
 \vspace*{2.6cm}
 {\raggedright\sffamily\bfseries\color{DNAblue}\Large PART \thepart\par}
 \vspace{4pt}
 {\raggedright\sffamily\bfseries\color{Ink}\fontsize{26}{31}\selectfont #1\par}
 \vspace{8pt}{\color{DNAblue}\rule{\linewidth}{1.2pt}}\par\vspace{10pt}
 {\raggedright\setlength{\parskip}{6pt}#3\par}%
 \clearpage}
\renewcommand{\thepart}{\Roman{part}}
\renewcommand{\cftpartfont}{\sffamily\bfseries\color{DNAblue}}
\renewcommand{\cftpartpagefont}{\sffamily\bfseries\color{DNAblue}}
